\selectlanguage{russian}

Большие языковые модели (LLM) продемонстрировали выдающиеся способности в решении разнообразных задач, однако их внедрение сталкивается с тремя критическими проблемами: ненадёжные выходные данные, которые могут содержать ошибки или галлюцинации, ограниченный контроль над генерацией и неэффективные процедуры обучения, расходующие вычислительные ресурсы на избыточные данные. Данная диссертация решает эти проблемы с помощью единого подхода, использующего сигналы прямого прохода модели --- как внутренние представления, так и выходные распределения --- для повышения надёжности, управляемости и эффективности обучения.

Мы представляем три взаимосвязанных исследования. Во-первых, мы исследуем оценку неопределённости путём анализа связи между энтропией на уровне токенов и правильностью ответов. Наши эксперименты на MMLU-Pro с четырьмя моделями (Phi-4, Mistral-Small-24B, Qwen-1.5B, Qwen-72B) показывают, что энтропия ответного токена достигает ROC-AUC до 0.83 для предсказания неправильных ответов в доменах, зависящих от знаний, тогда как подходы Model-as-Judge работают почти случайно. Во-вторых, мы исследуем геометрическую структуру мультиязычных представлений в латентном пространстве LLM, обнаруживая, что языковые направления могут быть идентифицированы с помощью PCA с точностью классификации 96--100\% для четырёх языковых пар. Проецируя компоненту исходного языка из скрытых состояний, мы уменьшаем дистрибуционную дивергенцию до 55\% и достигаем снижения индекса переключения кода на 63--95\% в сгенерированном тексте без дополнительного обучения. В-третьих, мы разрабатываем конвейер дообучения с учётом сложности, который использует энтропию для стратификации обучающих данных и применяет целевые интервенции: стандартное дообучение для обычных примеров и дистилляцию цепочки рассуждений для сложных. Этот подход повышает точность с 0.39 до 0.52 (Qwen-3B) и с 0.51 до 0.64 (Phi-4-Mini) на MMLU-Pro при использовании на 79--82\% меньше обучающих данных, с обобщением результатов на датасеты GSM8K и MedMCQA.

В совокупности эти вклады демонстрируют, что доступные сигналы одного прямого прохода предоставляют ценную информацию для решения фундаментальных проблем внедрения LLM. Разработанные методы не требуют архитектурных модификаций или дорогостоящего переобучения, что делает их практичными для реальных приложений.
